\chapter{门罗币金额隐藏 (Monero Amount Hiding)}
\label{chapter:pedersen-commitments}

在大多数像比特币这样的加密货币中，交易 output 票据——赋予花费一定"金额"资金的权利——以明文传达这些金额。这使得观察者可以轻松验证花费的金额等于发送的金额。

在门罗币中，我们使用{承诺 (commitment)}来对除发送者和接收者以外的所有人隐藏 output 金额，同时仍然让观察者确信交易发送的金额不多也不少。正如我们将看到的，金额承诺还必须有相应的"范围证明 (range proof)"来证明隐藏的金额在合法范围内。



\section{承诺 (Commitments)}
\label{sec:commitments}

一般来说，密码学中的{承诺方案 (commitment scheme)}是一种在不揭示值本身的情况下对一个值做出承诺的方式。一旦对某事做出承诺，你就必须坚守它。

例如，在一个抛硬币游戏中，Alice 可以通过将她承诺的值与秘密数据进行 hash 并发布 hash 值来私下承诺一个结果（即"猜正反"）。在 Bob 抛过硬币之后，Alice 宣布她承诺了哪个值，并通过揭示秘密数据来证明。然后 Bob 可以验证她的声明。

换句话说，假设 Alice 有一个秘密字符串 $blah$，她想要承诺的值是 $heads$。她计算 $h = \mathcal{H}(blah, heads)$ 并将 $h$ 给 Bob。Bob 抛硬币，然后 Alice 告诉 Bob 秘密字符串 $blah$ 并且她承诺了 $heads$。Bob 计算 $h' = \mathcal{H}(blah, heads)$。如果 $h' = h$，那么他就知道 Alice 在抛硬币之前就猜了 $heads$。

Alice 使用所谓的"盐 (salt)" $blah$，这样 Bob 就不能在抛硬币之前直接猜 $\mathcal{H}(heads)$ 和 $\mathcal{H}(tails)$，然后推断出她承诺了 $heads$。\footnote{如果承诺的值很难猜测和验证，例如如果它是一个看似随机的椭圆曲线点，那么给承诺加盐就不是必要的。}


\section{Pedersen commitment}
\label{pedersen_section}

{Pedersen commitment} \cite{Pedersen1992} 是一种具有{加法同态 (additively homomorphic)}性质的承诺。如果 \(C(a)\) 和 \(C(b)\) 分别表示值 \(a\) 和 \(b\) 的承诺，那么 \(C(a + b) = C(a) + C(b)\)。\footnote{在这个上下文中，加法同态意味着当你将标量通过标量 $x$ 的映射 $x \rightarrow x G$ 变换为 EC 点时，加法得以保持。}这个性质在承诺交易金额时非常有用，因为人们可以证明，例如，input 等于 output，而无需揭示手头的金额。\\

幸运的是，Pedersen commitment 很容易用椭圆曲线密码学实现，因为以下等式显然成立：\[a G + b G = (a + b) G\]

显然，如果将承诺简单地定义为 \(C(a) = a G\)，我们可以很容易地创建承诺的作弊表来帮助我们识别 $a$ 的常见值。

要获得信息论隐私，需要添加一个秘密{盲因子 (blinding factor)}和另一个生成元 \(H\)，使得对于哪个 \(\gamma\) 值以下等式成立是未知的：\(H = \gamma G\)。离散对数问题的困难性确保了从 $H$ 计算 $\gamma$ 是不可行的。\footnote{在门罗币的情况下\marginnote{src/ringct/ rctTypes.h}，$H = 8*to\_point(\mathcal{H}_n(G))$。这与 $\mathcal{H}_p$ hash 函数的不同之处在于，它直接将 $\mathcal{H}_n(G)$ 的输出解释为压缩点坐标，而不是通过数学方法派生曲线点（参见 \cite{hashtopoint-writeup}）。这种差异的历史原因\marginnote{tests/unit\_ tests/ ringct.cpp {\tt TEST(ringct, HPow2)}}我们不得而知，实际上这是唯一不使用 $\mathcal{H}_p$ 的地方（Bulletproofs 也使用 $\mathcal{H}_p$）。注意有一个 $*8$ 操作，以确保结果点在我们的 $l$ 子群中（$\mathcal{H}_p$ 也这样做）。}

%在门罗币的情况下\marginnote{src/ringct/ rctTypes.h}，$H = \mathcal{H}_p(G)$。\footnote{\label{hashtopoint_note}门罗币独有的 hash to point 函数\marginnote{src/ringct/ rctOps.cpp {\tt hash\_to\_p3()}}（参见 \cite{hashtopoint-writeup}）在实践中用于将一个曲线点的普通 hash 直接转换为另一个曲线点。对于承诺，$H = hash\_to\_point(\mathit{Keccak}(G))$。}%see rctTypes.h

然后我们可以将对值 \(a\) 的承诺定义为 \(C(x, a) = x G + a H\)，其中 \(x\) 是盲因子（又名"掩码 (mask)"），用于防止观察者猜测 $a$。

承诺 $C(x, a)$ 在信息论上是私密的，因为有许多可能的 $x$ 和 $a$ 组合会输出相同的 $C$。\footnote{基本上，有许多 $x'$ 和 $a'$ 使得 $x'+a' \gamma = x+a \gamma$。承诺者知道一种组合，但攻击者无法知道是哪一种。这种性质也被称为"完美隐藏 (perfect hiding)" \cite{adam-zero-to-bulletproofs}。此外，即使是承诺者也不能在不解决 $\gamma$ 的离散对数问题的情况下找到另一种组合，这种性质称为"计算绑定 (computational binding)" \cite{adam-zero-to-bulletproofs}。}如果 $x$ 是真正随机的，攻击者将完全没有方法推断出 $a$ \cite{maxwell-ct, SCOZZAFAVA1993313}。%{https://people.xiph.org/~greg/confidential_values.txt}


\section{金额承诺 (Amount commitments)}
\label{sec:pedersen_monero}

在门罗币中，output 金额以 Pedersen commitment 的形式存储在交易中。我们将对 output 金额 $b$ 的承诺定义为：\vspace{.175cm}
\[C(y,b) = y G + b H\marginnote{src/ringct/ rctOps.cpp {\tt addKeys2()}}\]

接收者应该能够知道他们拥有的每个 output 中有多少钱，以及重建金额承诺，以便它们可以作为新交易的 input。这意味着盲因子 $y$ 和金额 $b$ 必须传达给接收者。

采用的解决方案是使用 Diffie-Hellman 共享密钥 $r K_B^v$（回顾第 \ref{sec:multi_out_transactions} 节中的"交易公钥"）。对于区块链中的任何给定交易，其每个 output $t \in \{0, ..., p-1\}$ 都有一个掩码 $y_t$（发送者和接收者可以私下计算）和一个存储在交易数据中的{金额}。虽然 $y_t$ 是椭圆曲线标量，占 32 字节，但 $b$ 将被范围证明限制为 8 字节，因此只需要存储一个 8 字节的值。\footnote{与\marginnote{src/crypto- note\_core/ cryptonote\_ tx\_utils.cpp {\tt construct\_ tx\_with\_ tx\_key()} 调用 {\tt generate\_ output\_ ephemeral\_ keys()}}第 \ref{sec:one-time-addresses} 节中的一次性地址 $K^o$ 一样，output 索引 $t$ 在 hash 前被附加到共享密钥中。这确保了发往同一地址的 output 不会有相似的掩码和{金额}，概率可忽略不计。也和之前一样，项 $r K^v_B$ 乘以 8，所以实际上是 $8rK^v_B.$}\footnote{这个解决方案（在协议 v10 中实现）替代了之前使用更多数据的方法，从而导致交易类型从 v3（{\tt RCTTypeBulletproof}）变为 v4（{\tt RCTTypeBulletproof2}）。本报告的第一版讨论了之前的方法 \cite{ztm-1}。}\marginnote{src/ringct/ rctOps.cpp {\tt ecdh- Encode()}}[.725cm]\vspace{.175cm}%Under construct_tx_with_tx_key, key_amounts is created with the shared secret r K_B^v concatenated with the index by calling generate_output_ephemeral_keys() which then uses derivation_to_scalar(), then ecdhEncode builds the mask and amount when key_amounts is passed to it with the unmasked mask & amount. After update, ecdhHash and xor8 also play a role.
\begin{align*}
  y_t &= \mathcal{H}_n(``commitment\_mask",\mathcal{H}_n(r K_B^v, t)) \\
  \mathit{amount}_t &= b_t \oplus_8 \mathcal{H}_n(``amount", \mathcal{H}_n(r K_B^v, t))
\end{align*}

这里，$\oplus_8$ 表示对每个操作数的前 8 字节执行 XOR 操作（第 \ref{sec:XOR_section} 节）（$b_t$ 已经是 8 字节，$\mathcal{H}_n(...)$ 是 32 字节）。接收者可以对 $\mathit{amount}_t$ 执行相同的 XOR 操作来揭示 $b_t$。

接收者 Bob 将能够使用交易公钥 $r G$ 和他的查看密钥 $k_B^v$ 计算盲因子 $y_t$ 和金额 $b_t$。他还可以检查交易数据中提供的承诺 $C(y_t, b_t)$（此后记为 $C_t^b$）是否与手头的金额一致。\\

更一般地说，任何可以访问 Bob 查看密钥的第三方都可以解密他的 output 金额，还可以确保它们与关联的承诺一致。



\section{RingCT 简介 (RingCT introduction)}
\label{sec:ringct-introduction}

一笔交易将包含对其他交易 output 的引用（告诉观察者哪些旧的 output 将被花费），以及它自己的 output。output 的内容包括一个一次性地址（分配 output 的所有权）和一个隐藏金额的 output 承诺（以及第 \ref{sec:pedersen_monero} 节中的编码 output 金额）。

虽然交易的验证者不知道每个 input 和 output 中包含多少钱，但他们仍然需要确保 input 金额之和等于 output 金额之和。门罗币使用一种叫做 RingCT \cite{MRL-0005-ringct} 的技术来实现这一点，该技术于 2017 年 1 月首次实现（协议 v4）。

如果我们有一笔包含金额 \(a_1, ..., a_m\) 的 $m$ 个 input 的交易，以及金额 \(b_0, ..., b_{p-1}\) 的 $p$ 个 output，那么观察者有理由期望：\footnote{如果预期的 output 总金额不精确等于所拥有的 output 的任何组合，那么交易作者可以添加一个"找零" output 将多余的钱发回给自己。打个比方，用一张 20 美元的钞票支付 15 美元的开销，你会从收银员那里收到 5 美元。}\vspace{.175cm}
\[\sum_j a_j - \sum_t b_t = 0\]

由于承诺是加法的，我们不知道 $\gamma$，我们可以轻松地通过使 input 和 output 金额的承诺之和为零（即通过设置 output 盲因子之和等于旧 input 盲因子之和）来向观察者证明我们的 input 等于 output：\footnote{回忆第 \ref{elliptic_curves_section} 节，我们可以通过反转坐标然后相加来减去一个点。如果 $P = (x, y)$，$-P = (-x, y)$。还请回忆，域元素的取反是计算$\\pmod q$，所以 $(–x \\pmod q)$。}\vspace{.175cm}
\[\sum_{j}{C_{j, in}} - \sum_{t}{C_{t, out}} = 0\]

为了避免发送者可识别性，我们使用一种略有不同的方法。被花费的金额对应于之前交易的 output，这些 output 有承诺\vspace{.175cm}
\[C^a_{j} = x_j G + a_j H\]

发送者可以创建对相同金额但使用不同盲因子的新承诺；即，
\[C'^a_{j} = x'_j G + a_j H\]

显然，她会知道两个承诺之差的私钥：\vspace{.175cm}
\[C^a_{j} - C'^a_{j} = (x_j - x'_j) G\]\\
因此，她将能够使用这个值作为{零承诺 (commitment to zero)}，因为私钥 $(x_j - x'_j) = z_j$ 可以用来做签名，从而证明和中没有 $H$ 分量（假设 $\gamma$ 是未知的）。换句话说，证明 $C^a_{j} - C'^a_{j} = z_j G + 0H$，我们在讨论 RingCT 交易结构时将在第 \ref{chapter:transactions} 章中实际执行这一操作。

让我们称 $C'^a_j$ 为{伪 output 承诺 (pseudo output commitment)}。伪 output 承诺包含在交易数据中，每个 input 一个。

在将交易提交到区块链之前，网络将想要验证其金额是平衡的。伪承诺和 output 承诺的盲因子被选择为使得\vspace{.175cm}
\[\sum_j x'_j  - \sum_t y_t = 0\]

这\marginnote{src/ringct/ rctSigs.cpp verRct- Semantics- Simple()}允许我们证明 input 金额等于 output 金额：\vspace{.175cm}
\[(\sum_j C'^a_{j} - \sum_t C^b_{t}) = 0\]

幸运的是，选择这样的盲因子很容易。在当前版本的门罗币中，除了第 $m$ 个伪 output 承诺外，所有盲因子都是随机的，其中 $x'_m$ 简单地是\marginnote{genRct- Simple()}
\[x'_m = \sum_t y_t - \sum_{j=1}^{m-1} x'_j\]



\section{范围证明 (Range proofs)}
\label{sec:range_proofs}

加法承诺的一个问题是，如果我们有承诺 $C(a_1)$、$C(a_2)$、$C(b_1)$ 和 $C(b_2)$，并且我们打算用它们来证明 $(a_1 + a_2) - (b_1 + b_2) = 0$，那么如果等式中有一个值是"负的"，这些承诺仍然成立。

例如，我们可以有 $a_1 = 6$、$a_2 = 5$、$b_1 = 21$ 和 $b_2 = -10$。\vspace{.175cm}
\begin{flalign*}
    && (6 + 5) - (&21 + -10) = 0&\\
     \intertext{\quad \quad \quad \quad \quad where} && 21G + -10G = 21G + (&l-10)G = (l + 11)G = 11G&
\end{flalign*}

由于 $-10 = l-10$，我们实际上创造了比我们投入的多 $l$ 个门罗币（超过 7.2x10$^{74}$ 个！）。

%Bulletproofs start here
解决门罗币中这个问题的方案是使用 Benedikt B\"{u}nz 等人在 \cite{Bulletproofs_paper} 中首次描述的 Bulletproofs 证明方案（也参见 \cite{adam-zero-to-bulletproofs,dalek-bulletproofs-notes} 中的解释）来证明每个 output 金额在某个范围内（从 0 到 $2^{64}-1$）。\footnote{可以想象，如果几个 output 都在合法范围内，它们金额之和可能会溢出并导致类似的问题。然而，当最大 output 远小于 $l$ 时，需要大量的 output 才会发生这种情况。例如，如果范围是 0-5，l = 99，那么要使用一个输入 2 来伪造货币，我们需要 $5 + 5 + …. + 5 + 1 = 101 \equiv 2 \pmod{99}$，总共需要 21 个 output。在门罗币中 $l$ 大约比可用范围大 2\^{}189 倍，这意味着需要荒谬的 2\^{}189 个 output 才能伪造货币。}鉴于 Bulletproofs 的复杂和精密性质，本文档不对其进行阐述。此外，我们认为引用的材料充分地呈现了其概念。\footnote{在协议 v8 之前，范围证明是通过 Borromean 环签名完成的，在 Zero to Monero 的第一版中有所解释 \cite{ztm-1}。}

Bulletproof 证明算法\marginnote{src/ringct/ rctSigs.cpp {\tt proveRange- Bullet- proof()}}以 output 金额 $b_t$ 和承诺掩码 $y_t$ 作为输入，输出所有 $C^b_t$ 和一个 $n$ 元组聚合证明 $\Pi_{BP} = (A, S, T_1, T_2, \tau_x, \mu, \mathbb{L}, \mathbb{R}, a, b, t)$\footnote{向量 $\mathbb{L}$ 和 $\mathbb{R}$ 各包含 $\lceil \textrm{log}_2(64 \cdot p) \rceil$ 个元素。$\lceil$ $\rceil$ 表示 log 函数向上取整。由于它们的构造方式，某些 Bulletproofs 使用"虚拟 output"作为填充，以确保 $p$ 加上虚拟 output 的数量是 2 的幂。这些虚拟 output 可以在验证期间生成，不与证明数据一起存储。}\footnote{Bulletproof 中的变量与本文档中的其他变量无关。符号重叠纯属巧合。注意群元素 $A, S, T_1, T_2, \mathbb{L},$ 和 $\mathbb{R}$ 在存储前乘以 1/8，在验证期间乘以 8。这确保它们都是 $l$ 子群的成员（回顾第 \ref{elliptic_curves_section} 节）。}。该单一证明用于同时证明所有 output 金额在范围内，聚合它们大大减少了空间需求（尽管确实增加了验证时间）。\footnote{实际上，多个单独的 Bulletproofs 可以被一起"批处理"\marginnote{rct/ringct/ bulletproofs.cc {\tt bulletproof\_ VERIFY()}}，这意味着它们被同时验证。这样做改善了验证它们所需的时间，目前在门罗币中 Bulletproofs 是按区块批处理的，尽管理论上批处理多少没有限制。每笔交易只允许有一个 Bulletproof。}验证算法以所有 $C^b_t$ 和 $\Pi_{BP}$ 作为输入，如果所有承诺的金额都在 0 到 $2^{64} - 1$ 的范围内则输出 {\tt true}。

$n$ 元组 $\Pi_{BP}$ 占用 $(2 \cdot \lceil \textrm{log}_2(64 \cdot p) \rceil + 9) \cdot 32$ 字节的存储空间。